\documentclass[12pt]{article}
\usepackage[utf8]{inputenc}
\usepackage[T1]{fontenc}
\usepackage{hyperref}
\title{Congo Swahili Language Resources for K-12 Curriculum Development}
\author{Compiled by Claude Sonnet 5.5 (claude-sonnet-5-5) for Andr\'as Kornai}
\date{}
\begin{document}
\maketitle

\section{Introduction}
Congo Swahili covers the Swahili varieties of the Democratic Republic of the Congo (Katanga, Kivu, Lualaba, Kingwana and others). Swahili is one of the DRC's four national languages, and French is the official language. Its role in schools differs by region. Few dedicated educational resources turned up in these searches.

\section{Primary Educational Resources}
No Congo Swahili primers or textbooks were found in these searches.

\section{Linguistic and Reference Materials}
\begin{itemize}
\item Wikipedia overview of Congo Swahili, which cites W.\ F.\ Hunter's 1959 \textit{A Manual of Congo Swahili Grammar}: \url{https://en.wikipedia.org/wiki/Congo_Swahili}
\item Glottolog reference entry for Hunter (1959): \url{https://glottolog.org/resource/reference/id/154276}
\item Linguistic study of Lubumbashi Swahili (dissertation attachment on Qucosa): \url{https://ul.qucosa.de/api/qucosa%3A14609/attachment/ATT-0/}
\item Elon grammar page on Swahili in the DRC: \url{https://elon.io/grammar/swahili/countries/democratic-republic-of-congo}. The search results did not assess its quality.
\end{itemize}

\section{Religious and Community Resources}
\begin{itemize}
\item Bible.is entries for Swahili (\texttt{swc}) via ScriptureEarth: \url{https://www.scriptureearth.org/00i-Scripture_Index.php?iso=swc}
\end{itemize}

\section{Digital and Audio Resources}
The only audio found is the Bible.is material linked above.

\section{Curriculum Development Context}
Little dedicated material for Congo Swahili exists in these results. Standard Swahili materials may serve as a base, but the Congolese varieties differ: they take many French loans, have a different noun-class system, and one source reports a proposed orthography using ``c'' for ``ch''. The Wikipedia figure of about 2 million first-language speakers is much lower than the survey's reference figure.

\end{document}